\documentclass{ximera}

\title{Optimization and Related Rates Review Activity Solution Guide}
\author{MATH 425: Calculus I}

\begin{document}
\begin{abstract}
   Working with the peers in your group, solve the following problems. Make sure to show and justify all your work. Make sure everyone in the group understands the solution and participates. Be prepared to report your answers to the whole class. 
\end{abstract}
\maketitle

This activity will help you review and understand the process for optimization and related rates in context.

\begin{exercise}
Suppose contestants with potentially different running speeds are planning for an unusual running race.  In this race, a relay team consists of exactly two runners, and those runners must travel on diagonal paths in a zig-zag shape.  The runners are able to choose where they hand off the relay baton, so each team is interested in planning a strategy which \textbf{minimizes} their combined travel time. \\
Let us suppose the race course is 80 m long, and 20 m wide. The first contestant starts at the corner of the race course and runs to the other side at whatever angle they choose.  When they reach the opposite side of the race course, they hand off to the second contestant who runs towards the far corner of the course.\\
The objective function for travel time for a single relay team is constructed as follows: $$T(x)=\frac{zig(x)}{s_1}+\frac{zag(x)}{s_2}$$
Where $zig(x)$ and $zag(x)$ are the distances travelled by the first and second runners (respectively), and the $s_i$ are their respective speeds.\\
Use the geometry of the race course and the Pythagorean Theorem, assuming that the distance from the opposite side corner to the place where the handoff takes place is $x$, find an objective function in terms of $x$.  If the runners happen to run exactly at the same speed (they trained together for this race, after all), where should the hand off occur? (Namely, what value of $x$ results in the minimum travel time?)

Where should the contestants plan to hand over the baton?\\
\textcolor{blue}{
    \begin{image}
        \includegraphics{opitimization race.png}
    \end{image}
    \xmalt{A rectangle with the top left and right hand corners labelled ``start'' and ``finish'', respectively.  The top left corner ``start'' point is connected by a line to a point on the bottom of the rectangle labelled ``hand off point.'' This point is also connected by a line to the top right corner ``finish'' point.}
We assume $x$ is the distance to the left of the hand off point, making the distance to the right of the hand of point $80-x$.\\
If $zig(x)$ is the distance run by the first runner, $zig(x)=\sqrt{400+x^2}$ by the Pythagorean Theorem.\\
If $zag(x)$ is the distance run by the second runner, $zag(x)=\sqrt{400+(80-x)^2}$ by the Pythagorean Theorem.\\
$$T(x)=\frac{\sqrt{400+x^2}}{s}+\frac{\sqrt{400+(80-x)^2}}{s}$$
$$T'(x)=\frac{1}{s}\bigg(\frac12(400+x^2)^{-1/2}(2x)+\frac12(400-(80-x)^2)^{-1/2}(2)(80-x)(-1)\bigg)=\frac{1}{s}\bigg(\frac{x}{\sqrt{400+x^2}}-\frac{80-x}{\sqrt{400+(80-x)^2}}\bigg)$$\\
To find the minimum value of $T(x)$, we need to check where $T'(x)=0$ and where $T'(x)$ DNE. \\
$\dfrac{1}{s}\bigg(\dfrac{x}{\sqrt{400+x^2}}-\dfrac{80-x}{\sqrt{400+(80-x)^2}}\bigg)=0$ when $\bigg(\dfrac{x}{\sqrt{400+x^2}}-\dfrac{80-x}{\sqrt{400+(80-x)^2}}\bigg)=0$\\
Or, when $\dfrac{x}{\sqrt{400+x^2}}=\dfrac{80-x}{\sqrt{400+(80-x)^2}}$.  }\\
\textcolor{blue}{
Multiplying by both denominators gives: \\
$x\sqrt{400+(80-x)^2}=(80-x)\sqrt{400+x^2}$\\
Square both sides: \\
$x^2(400+(80-x)^2)=(80-x)^2(400+x^2)$\\
Distribute: \\
$400x^2+x^2(80-x)^2=400(80-x)^2+x^2(80-x)^2$\\
Subtract the common term from both sides:\\
$400x^2=400(80-x)^2$\\
Divide by 400 and expand $(80-x)^2$:\\
$x^2=6400-160x+x^2$\\
$6400-160x=0$, or $x=40$m.  \\
%
Check when $T'$ DNE: both denominators are distances, so they will always be positive values, and thus the expressions will always exist.\\
%
Classify the absolute minimum:
Domain of $T(x)$: $[0,80]$\\
$T(0)=\frac{1}{s}\bigg(\sqrt{400+0^2}+\sqrt{400+80^2}\bigg)=\frac{102.46}{s}$\\
$T(40)=\frac{1}{s}\bigg(\sqrt{400+40^2}+\sqrt{400+40^2}\bigg)=\frac{89.44}{s}$\\
$T(80)=\frac{1}{s}\bigg(\sqrt{400+80^2}+\sqrt{400+0^2}\bigg)=\frac{102.46}{s}$\\
%
$T(40)$ is the smallest, making $x=40$ the distance giving the shortest time.}
\end{exercise}

\begin{exercise}
Suppose you are video recording the ascent of a hot air balloon.  The balloon is rising from a spot on the ground 300 ft away from your location, at a constant rate of 1 ft/sec.  When the camera is at an angle of $\frac{\pi}{6}$ from horizontal, what is the rate of change of the angle of the camera tracking the balloon?  Assume the camera is at ground level.\\
\textcolor{blue}{
    \begin{image}
        \includegraphics{RelatedRatesReview.png}
    \end{image}
    \xmalt{A triangle on a coordinate plane with base of length 300, labelled $c$ on the $x$-axis, height of 400 labelled $a$ parallel to the $y$-axis on the line $x=300$, and hypotenuse labelled $b$ connecting the endpoints of the base and height. Each angle of the triangle is labelled with the captial letter matching the side opposite it.}
Before we solve the problem, let's think about all possibilities for equations for this triangle: \\
$b^2=a^2+c^2$\\
$\tan(A)=\frac{a}{c}$\\
$\sin(A)=\frac{a}{b}$\\
$\cos(A)=\frac{c}{b}$\\
Which makes the most sense to use?  List the given information:\\
$c=300$\\
$a$ not given\\
$\frac{dc}{dt}=0$\\
$\frac{da}{dt}=1$\\
$A=\frac{\pi}{6}$\\
$\frac{dA}{dt}=$unknown}\\
\textcolor{blue}{
%
What equation makes the most sense?  The most given and unknowns are involved in $\tan(A)=\frac{a}{c}$.\\
Take the derivative with respect to $t$: $\sec^2(A)\frac{dA}{dt}=\frac{1}{c}\frac{da}{dt}$\\
Use knowns:\\
$$\sec^2\bigg(\frac{\pi}{6}\bigg)\frac{dA}{dt}=\frac{1}{300}(1)$$\\
%
$\sec(\frac{\pi}{6})=\frac{1}{\cos(A)}=\frac{2}{\sqrt{3}}$\\
So, $$\bigg(\frac{2}{\sqrt{3}}\bigg)^2\frac{dA}{dt}=\frac{1}{300}$$
$$\frac{dA}{dt}=\frac{1}{300}\bigg(\frac{3}{4}\bigg)=\frac{1}{400}$$ radians/sec}
\end{exercise}



\end{document}
